\section{The counterexample and its localizations}
\label{scope:section}

The algebraic construction and obstruction now give all three inputs to
Proposition~\ref{transfer:criterion}.  We first determine the exact
fp-type of its connective spectrum and prove both localization comparisons
for that same spectrum.  These conclusions finish the main theorem.
We then explain why its positive local classification is compatible with
the failure of global thick generation, and isolate the separate obstruction
for the connective cover of Morava \(E\)-theory.

\subsection{Exact type and telescope comparisons}

Fix a prime satisfying the operator bounds in \eqref{op:parameters}.  Let
\(U\) be the object of Proposition~\ref{glue:object}, and write
\[
 Z_{\rm sp}\in\Sp_{E(2)},\qquad
 X_{\rm sp}=\tau_{\geq0}Z_{\rm sp}
\]
for the spectra obtained from \(U\) by Proposition~\ref{transfer:criterion}.
That proposition gives finite \(p\)-groups \(\pi_iZ_{\rm sp}\) in every
degree, identifies \(L_{E(2)}X_{\rm sp}\) with \(Z_{\rm sp}\), and shows
that \(X_{\rm sp}\) is \(p\)-complete of fp-type at most two with
\[
 X_{\rm sp}\notin\Thick(b_p^\wedge).
\]
Here \(b\) is the regular \(\BP\)-module quotient representing
\(\BP\langle2\rangle\) in Proposition~\ref{comp:generators}.  We first
strengthen the upper bound on fp-type to an equality.

\begin{proposition}
\label{scope:exact-height}
The spectrum \(X_{\rm sp}\) has fp-type exactly two.
\end{proposition}

\begin{proof}
Suppose that \(X_{\rm sp}\) had fp-type at most one.  The height-one
theorem of Lee--Pstr{\k a}gowski
\cite[Theorems~1.3 and~5.4]{LP2026} would give
\[
 X_{\rm sp}\in\Thick(\BP\langle1\rangle_p^\wedge).
\]
We check that the generator on the right is already in
\(\Thick(b_p^\wedge)\).

In the same monoidal category of \(\BP\)-modules used to construct
\(b\), form the further scalar cofiber
\[
 b_1=\cofib\bigl(\Sigma^{2(p^2-1)}b
                  \xrightarrow{v_2}b\bigr).
\]
Multiplication by \(v_2\) is regular on
\(\pi_*b=\Z_{(p)}[v_1,v_2]\), so
\(\pi_*b_1=\Z_{(p)}[v_1]\).  The mod-\(p\) comparison in
Proposition~\ref{comp:standard-cohomology}, with the cutoff changed from two to
one, identifies
\[
 H^*(b_1;\F_p)
   \cong\mathcal A//E(Q_0,Q_1).
\]
Indeed the same Koszul calculation has exterior generators
\(\epsilon_i\) for \(i>1\); its comparison with the calculation for
\(H\F_p\) is injective, so the spectral sequence collapses and the
Milnor-basis dimension comparison identifies the displayed Steenrod
module.  Here \(E(Q_0,Q_1)\) is the exterior subalgebra generated by the
first two Milnor primitives.  The quotient is connective and of finite
type over \(\Z_{(p)}\).  Lee's odd-prime uniqueness theorem therefore
identifies \((b_1)_p^\wedge\) with
\(\BP\langle1\rangle_p^\wedge\) as a spectrum
\cite[Theorem~1.1]{Lee2026}; see also \cite[Note~2]{LP2026}.

Since \(p\)-completion is exact, completing the defining cofiber
sequence of \(b_1\) gives
\[
 \BP\langle1\rangle_p^\wedge\in\Thick(b_p^\wedge).
\]
The assumed height-one membership would consequently put
\(X_{\rm sp}\) in \(\Thick(b_p^\wedge)\), contradicting the transfer
criterion.  Its fp-type is at most two and not at most one, as required.
\end{proof}

Write \(L_2=L_{E(2)}\), and let \(L_2^f\) denote finite chromatic
localization away from finite spectra of type greater than two.  For
\(j>0\), let \(T(j)\) be the telescope of a positive-degree
\(v_j\)-self-map on a finite type-\(j\) spectrum, and put
\(T(0)=H\Q\).  These localizations occur in the restricted telescope
conjecture of Mahowald--Rezk
\cite[Conjecture~7.3]{MR1999}, and in the proposed reconstruction route
of \cite[Conjecture~1.2 and Remark~5.1]{LP2026}.  The following
comparisons are proved directly for the spectrum constructed here.

\begin{proposition}
\label{scope:telescopes}
There are equivalences
\[
 L_2^fX_{\rm sp}\simeq L_2X_{\rm sp}\simeq Z_{\rm sp},
 \qquad
 L_{T(2)}X_{\rm sp}\simeq L_{K(2)}X_{\rm sp}.
\]
\end{proposition}

\begin{proof}
Put \(N=\tau_{<0}Z_{\rm sp}\), so that
\[
 X_{\rm sp}\longrightarrow Z_{\rm sp}\longrightarrow N
\]
is a cofiber sequence.  The spectrum \(N\) is the filtered colimit of
its finite Postnikov ranges \(\tau_{\geq-m}N\).  Every homotopy group
of each such range is a finite \(p\)-group.  Filtering those groups by
their composition series expresses each range as a finite extension of
shifts of \(H\F_p\).

For \(j>0\), the ordinary mod-\(p\) homology of the finite spectrum
on which a \(v_j\)-self-map is chosen is bounded.  Iterating a
positive-degree self-map therefore gives the zero telescope on that
homology, and hence
\(H\F_p\wedge T(j)=0\).  The same conclusion for \(j=0\) is rational
acyclicity.  It follows that \(N\) is \(T(j)\)-acyclic for
\(0\leq j\leq2\).  Finite localization is smashing and its acyclics
are generated by a finite type-three spectrum
\cite[author-preprint Corollary~1.19(i) and Theorem~2.7(ii)]{Ravenel1993}.
Together with the Boolean algebra theorem
\cite[Theorem~7.2.9]{Ravenel1992}, this identifies
\(\langle L_2^fS^0\rangle\) with
\(\langle T(0)\vee T(1)\vee T(2)\rangle\).
Thus \(N\) is \(L_2^f\)-acyclic.

Every finite spectrum of type greater than two is \(E(2)\)-acyclic.
Thus an \(E(2)\)-local spectrum is also \(L_2^f\)-local: it is right
orthogonal to the localizing subcategory generated by those finite
spectra.  In particular \(Z_{\rm sp}\) is \(L_2^f\)-local, and the
preceding cofiber sequence proves
\(L_2^fX_{\rm sp}\simeq Z_{\rm sp}\).  The equality with
\(L_2X_{\rm sp}\) was established in the transfer criterion.  This
argument is consistent with the finite-localization property of the
connective-cover examples in \cite[Section~7.2(2)]{MR1999}.

For the second comparison we may work with \(Z_{\rm sp}\), since
\(N\) is \(T(2)\)-acyclic and hence also \(K(2)\)-acyclic.  The
latter implication follows from the field property: \(K(2)_*T(2)\)
is a nonzero graded vector space over \(K(2)_*\), and the K\"unneth
isomorphism shows that a \(T(2)\)-acyclic spectrum has zero
\(K(2)\)-homology.  The same implication shows that every
\(K(2)\)-local spectrum is \(T(2)\)-local.

Consider the fiber
\[
 G=\fib\bigl(Z_{\rm sp}\longrightarrow
                    L_{K(2)}Z_{\rm sp}\bigr).
\]
It is \(K(2)\)-acyclic and \(E(2)\)-local.  Here the target is
\(E(2)\)-local because every \(E(2)\)-acyclic spectrum is
\(K(2)\)-acyclic.  The Hopkins--Ravenel smashing theorem implies
that \(T(2)\wedge G\) remains \(E(2)\)-local
\cite[Theorem~7.5.6]{Ravenel1992}.  Its \(K(0)\)- and
\(K(1)\)-homology vanish because they vanish on \(T(2)\), while its
\(K(2)\)-homology vanishes because it vanishes on \(G\).  The
Bousfield-class identity
\[
 \langle E(2)\rangle
   =\langle K(0)\vee K(1)\vee K(2)\rangle
\]
therefore makes \(T(2)\wedge G\) also \(E(2)\)-acyclic
\cite[Theorem~7.3.2(d)]{Ravenel1992}.  Being both local and acyclic,
it is zero.  Thus the map from \(Z_{\rm sp}\) to its
\(K(2)\)-localization is a \(T(2)\)-equivalence with
\(T(2)\)-local target.  This proves the claimed equivalence of the
two localizations, and hence the assertion for \(X_{\rm sp}\).
\end{proof}

\begin{proof}[Proof of Theorem~\ref{thm:main}]
The parameter choices for Proposition~\ref{op:operator} are available
at every sufficiently large prime.  For any such prime,
Proposition~\ref{glue:object} provides an object \(U\) with finite
\(\pi_iU\) in every degree and
\(\widehat U\in\Thick(\widehat T')\).
Theorem~\ref{obs:nonmembership} gives
\(U\notin\Thick(T')\).  Applying
Proposition~\ref{transfer:criterion} produces the connective,
\(p\)-complete spectrum \(X_{\rm sp}\) of fp-type at most two
outside \(\Thick(\BP\langle2\rangle_p^\wedge)\).
Proposition~\ref{scope:exact-height} gives exact fp-type two, and
Proposition~\ref{scope:telescopes} gives both localization
comparisons.
\end{proof}

This gives the asserted failure at sufficiently large primes, without
a claim at the remaining primes.  The telescope comparisons are
properties of this example; they neither assume nor settle the general
restricted telescope conjecture for fp-spectra.

\subsection{The local classification and its hypotheses}

We now explain the positive local conclusion for this same example.
For general height \(n\), write \(E_n\) for Morava \(E\)-theory;
this notation is distinct from the Johnson--Wilson theory \(E(n)\).
A \(K(n)\)-local spectrum \(Y'\) is \emph{locally fp} if
\(\pi_i(F\wedge Y')\) is finite for every integer \(i\) and every
finite type-\(n\) spectrum \(F\).  This condition asks for
degreewise finiteness.  It does not ask that the direct sum over all
degrees be finite, as in the finite-spectrum test for fp-spectra.

\begin{corollary}
\label{scope:local-classification}
The spectrum \(Y'=L_{K(2)}X_{\rm sp}\) is locally fp.  In the
\(K(2)\)-local category it belongs to
\[
 \Thick(E_2)
   =\Thick\bigl(L_{K(2)}\BP\langle2\rangle\bigr).
\]
Its Morava \(E\)-cohomology \(E_2^*Y'\) is finitely generated over
the ring of operations \(E_2^*E_2\), and it is reflexive in the
\(K(2)\)-local category.
\end{corollary}

\begin{proof}
The exact fp-type from Proposition~\ref{scope:exact-height} and the
\(T(2)\)-to-\(K(2)\) comparison from
Proposition~\ref{scope:telescopes} verify the hypotheses of
\cite[Lemma~4.1]{LP2026}.  That lemma gives local fp-ness.  The
equivalent characterizations in \cite[Theorems~1.7, 1.8, and~4.3]{LP2026}
then give thick membership, finite generation over Morava operations,
and local reflexivity.
\end{proof}

By Corollary~\ref{scope:local-classification}, the spectrum
\(L_{K(2)}X_{\rm sp}\) satisfies the hypotheses of the local reduction
and resolution results discussed below.  Their operations are Morava
operations, and their finite constructions take place in the local
category.  They do not provide integral resolutions over global
\(\BP\langle2\rangle\)-operations or lifts of maps through the
height-open fracture square.  We spell out the relevant hypotheses,
since they specify exactly what finite constructions the local theorem
provides.

\begin{remark}[Periodic reductions in the local theorem]
\label{scope:local-reductions}
Theorem~3.9 of Lee--Pstr{\k a}gowski, stated as Theorem~1.9 in their
introduction, chooses a generalized type-\(n\) Moore spectrum \(M\)
and an ideal \(I\) generated by suitable powers of the generators of
the height ideal, so that
\[
 E_{n*}M\cong E_{n*}/I.
\]
The powers are chosen to permit the required multiplicative
structure.  They also choose a sufficiently small normal open subgroup
\(U_{\rm gp}\) of the extended Morava stabilizer group.  After
extension to the finite \(K(n)\)-local Galois extension
\(E_n^{hU_{\rm gp}}\), the reduction object
\[
 V'=E_n^{hU_{\rm gp}}\wedge M
\]
has homotopy
\[
 \pi_*V'
   \cong\bigwedge_{E_{n*}/I}
           (\bar x_1,\ldots,\bar x_{n^2}),
 \qquad |\bar x_i|=-1.
\]
Their lower \(p\)-series and trivial-action hypotheses ensure that
the subgroup acts trivially on the coefficient quotient, that its
cohomology has the required exterior form, and that the
Adams--Novikov spectral sequence collapses with the stated extensions.
The assertion concerns these choices of powers and subgroup; arbitrary
Moore exponents and arbitrary open subgroups need not have these
properties \cite[Section~3 and Theorem~3.9]{LP2026}.

The object \(V'\) is compact in the \(K(n)\)-local category and is
a retract of the localization of a finite type-\(n\) spectrum, as
used in \cite[Notation~4.6]{LP2026}.  Thus, if \(Y'\) is locally
fp, the groups
\(\pi_iL_{K(n)}(V'\wedge Y')\) are finite for all \(i\).  The
finite-smash version of this assertion needs no extra localization:
if \(F\) is finite and \(Y'\) is local, finite duality makes
\(F\wedge Y'\) local.

In fact \(V'\wedge Y'\) itself is \(K(n)\)-local here.  Both
\(E_n^{hU_{\rm gp}}\) and \(Y'\) are \(E(n)\)-local, so
\[
 A'=E_n^{hU_{\rm gp}}\wedge Y'
\]
is \(E(n)\)-local by smashing localization.  For a
\(K(n)\)-acyclic spectrum \(W\), finite duality gives
\[
 [W,M\wedge A']_*
     \cong[M^\vee\wedge W,A']_*.
\]
The spectrum on the left of the second mapping group is
\(E(n)\)-acyclic.  Its lower Morava \(K\)-theories vanish because
\(M^\vee\) has type \(n\), and its \(K(n)\)-homology vanishes
because that of \(W\) vanishes.  The mapping group is consequently
zero, proving the claim.  It follows that
\(\pi_*(V'\wedge Y')\) is degreewise finite.  Periodicity and the
Artin coefficient quotient then make it a finitely generated module
over the displayed graded exterior algebra.
\end{remark}

\begin{remark}[Operations and finite local resolutions]
\label{scope:local-resolutions}
The reductions of Remark~\ref{scope:local-reductions} are used in a
universal-coefficient spectral sequence for \(V'\)-linear cohomology
with target \(E_n\wedge M\).  This spectral sequence is a module
over the universal-coefficient spectral sequence for the
\(V'\)-linear endomorphisms of that target.  At the coefficient
level the relevant self-Ext algebra is
\[
 \Ext_{\pi_*V'}^*(E_{n*}/I,E_{n*}/I),
\]
a polynomial algebra on the exterior generators' Ext classes, all
even in total degree.  The even-degree endomorphism spectral sequence
collapses.  The accompanying noetherianity and Ext-finiteness
arguments give finite-page collapse and strong convergence for the
module spectral sequence, together with finite generation over
operations at the quotient level.  Derived completeness then permits
passage back from these reductions.  These are the steps in
\cite[Lemmas~4.7--4.9 and Proposition~4.11]{LP2026} that produce
finite generation of \(E_n^*Y'\) over \(E_n^*E_n\).

For the converse direction, the twisted group algebra description
gives noetherian operations and global dimension \(n^2+n\) for
suitable uniform open subgroups
\cite[Proposition~2.12 and Corollary~2.13]{LP2026}.  Taking the
subgroup to be pro-\(p\) and contained in the kernel of the Galois
quotient also makes finitely generated projectives free
\cite[Proposition~2.14 and Corollary~2.15]{LP2026}.  The finite
resolution and cofiber argument tracks projective dimension over this
small-subgroup algebra.  After a suitable
\(E_n^{hU_{\rm gp}}\)-extension it reduces to a free case and then
descends thick membership in the \(K(n)\)-local category
\cite[Lemmas~4.4--4.5 and the proof of Theorem~4.3]{LP2026}.

The subgroup filtrations and completed operations used in this paper
provide related control of the complete part of a gluing diagram.
After the field extension, characteristic reduction, and removal of
central inertia, the digit filtration, compact Koszul resolution, and
complete endomorphism calculation give the bounds and models proved
above.  The discrete lattice and the obstruction to global maps impose
additional conditions.  Theorem~\ref{obs:nonmembership} uses
those conditions to exclude a global finite construction even though
the corresponding local membership holds.
\end{remark}

\subsection{Connective covers and tensor-ideal classifications}

There is also a direct obstruction to carrying the height-one
connective-cover step of \cite[Lemma~5.3]{LP2026} over to Morava
\(E\)-theory at height two.

\begin{proposition}
\label{scope:connective-cover}
For height-two Morava \(E\)-theory,
\[
 \tau_{\geq0}E_2\notin\Thick(b_p^\wedge).
\]
\end{proposition}

\begin{proof}
For a spectrum \(P\), write \(P/p=\cofib(p:P\to P)\).  The
property that \(\pi_i(P/p)\) is finite for every \(i\) is thick:
the multiplication-by-\(p\) cofiber functor is exact, and
degreewise finiteness is preserved by shifts, finite sums, cofiber
sequences, and retracts.  The property holds for \(b_p^\wedge\),
since
\[
 \pi_*\bigl(b_p^\wedge/p\bigr)=\F_p[v_1,v_2]
\]
has finite-dimensional homogeneous pieces.

On the other hand,
\[
 \pi_0E_2=W(\F_{p^2})[[u_1]],
\]
where \(W(\F_{p^2})\) is the Witt-vector ring.  Since
\(\tau_{\geq0}E_2\) is connective, its multiplication-by-\(p\)
cofiber has
\[
 \pi_0\bigl((\tau_{\geq0}E_2)/p\bigr)
      \cong\F_{p^2}[[u_1]],
\]
which is infinite.  It cannot lie in the thick subcategory generated
by \(b_p^\wedge\).
\end{proof}

Consequently, taking connective covers of the pieces in a finite
\(K(2)\)-local construction and correcting by Postnikov sections
does not by itself reproduce the height-one argument.  In the
construction here, the connective cover is taken only after an
\(E(2)\)-local object with degreewise finite homotopy has been
glued.  The transfer criterion proves fp-ness for that connective
cover, while the global obstruction determines its failure of thick
generation.

\begin{remark}[Local tensor ideals and comodule realization]
\label{scope:tensor-ideals}
Barthel--Heard--Naumann study a different classification problem in
\emph{On conjectures of Hovey--Strickland and Chai}.  Their
Definition~3.2 and Conjecture~3.4 concern thick tensor ideals of
\(K(n)\)-locally dualizable spectra.  For \(0\leq j\leq n\),
their subcategory \(\mathcal D_j\) consists of retracts of
\(A\otimes F\), where \(A\) is locally dualizable, \(F\) is
finite of type at least \(j\), and the tensor product is taken in
the \(K(n)\)-local category; they set
\(\mathcal D_{n+1}=(0)\).  Theorem~4.9 deduces the proposed
classification from the Chai-type condition on invariant radical
ideals.  Theorem~4.13 gives the converse implication under
\(2p-2>n^2+n\), and Theorem~4.15 establishes the height-two
classification at every prime \cite{BHN2022}.

The corresponding local Hovey--Strickland classification is established
unconditionally for every prime and every height \(n\geq1\) in
\cite[Theorem~1.2]{OpenAIStabilizerIdeals2026}: the category of dualizable
\(K(n)\)-local spectra has exactly \(n+2\) thick tensor ideals.

The telescope-versus-Morava vanishing comparison used in Barthel--Heard--Naumann's
discussion of Proposition~3.3 also applies to our \(E(2)\)-local
calculation, as the proof of Proposition~\ref{scope:telescopes}
shows.  The classification itself has different hypotheses and
closure operations: its objects are locally dualizable, with the
corresponding finiteness of their Morava modules, and its subcategories
are closed under tensoring with locally dualizable objects.  The
ordinary global thick subcategory of fp-spectra considered in
Theorem~\ref{thm:main} is not such a tensor-ideal category.

A further point of contact is large-prime comodule realization.
\cite[Lemma~3.11]{BHN2022} recalls Pstr{\k a}gowski's realization
functor from \(E_{n*}E_n\)-comodules to the \(E_n\)-local homotopy
category under \(2p-2>n^2+n\).  Barthel--Heard--Naumann apply it to
invariant-ideal reductions and reassemble them to prove their
Theorem~3.7, realizing finitely generated Morava modules by local
dualizable spectra.  Our prime range satisfies this inequality at
\(n=2\).  However, the object constructed here includes a lattice
map and is glued in a derived category over heights at most two;
thick membership depends on maps between such objects, in addition
to the realization of their individual homology modules.

For this reason the transfer criterion uses the homology-compatible
triangulated comparison for \(E(2)\) in Theorem~\ref{comp:algebraicity},
together with the checks on pure representatives in
Proposition~\ref{comp:generators}; see
\cite[Corollary~6.9 in the cited preprint version]{Pstragowski2021}.
The comparison transports the derived object and the finite constructions
needed in the proof.  The reduced group action and orbit control for its
particular height-two charts were established directly in the preceding
algebraic sections.
\end{remark}
